% !TEX program = lualatex
% =============================================================
% OnePage de Evidencias - Equipe 3 (Hanaro / Scrap) . Sprint 02 (W36)
%   A3 paisagem, tema LG. Foco: Agentes ReAct (LangGraph) & Controle de Scrap.
%   Compilar: uv run build onepage-evidencias  (2 passadas).
% =============================================================
\documentclass[11pt]{article}
\usepackage[a3paper,landscape,top=0.8cm,bottom=1.0cm,left=1.25cm,right=1.25cm]{geometry}

\usepackage{fontspec}
\defaultfontfeatures{Ligatures=TeX}
\IfFileExists{../../fonts/LGEITextTTF-Regular.ttf}{%
  \setmainfont{LGEITextTTF-Regular}[Path=../../fonts/, Extension=.ttf, BoldFont=LGEITextTTF-Bold]
  \newfontfamily\lghead{LGEIHeadlineTTF-Bold}[Path=../../fonts/, Extension=.ttf]
}{%
  \IfFileExists{../fonts/LGEITextTTF-Regular.ttf}{%
    \setmainfont{LGEITextTTF-Regular}[Path=../fonts/, Extension=.ttf, BoldFont=LGEITextTTF-Bold]
    \newfontfamily\lghead{LGEIHeadlineTTF-Bold}[Path=../fonts/, Extension=.ttf]
  }{%
    \IfFileExists{fonts/LGEITextTTF-Regular.ttf}{%
      \setmainfont{LGEITextTTF-Regular}[Path=fonts/, Extension=.ttf, BoldFont=LGEITextTTF-Bold]
      \newfontfamily\lghead{LGEIHeadlineTTF-Bold}[Path=fonts/, Extension=.ttf]
    }{%
      \setmainfont{LG EI Text}[BoldFont=* Bold]
      \newfontfamily\lghead{LG EI Headline Bold}
    }%
  }%
}

\usepackage{xcolor}
\definecolor{lgred}{RGB}{165,0,52}
\definecolor{axdark}{RGB}{30,10,10}
\definecolor{cardline}{RGB}{206,170,172}
\definecolor{muted}{RGB}{110,95,95}
\definecolor{tintbg}{RGB}{250,242,244}

\usepackage{graphicx}
\usepackage{tikz}
\usetikzlibrary{calc}
\usepackage[most]{tcolorbox}
\usepackage{enumitem}
\usepackage{fontawesome5}
\usepackage[hidelinks]{hyperref}
\hypersetup{colorlinks=true, urlcolor=lgred}
\setlist[itemize]{leftmargin=1.15em, itemsep=2.5pt, topsep=1.5pt, parsep=0pt}

\pagestyle{empty}
\setlength{\parindent}{0pt}
\newcommand{\hd}[1]{{\lghead\color{lgred}#1}}

\newtcolorbox{card}[2][]{
  enhanced, colback=white, colframe=cardline, boxrule=0.6pt, arc=3pt,
  left=9pt, right=9pt, top=5pt, bottom=6pt,
  fonttitle=\lghead\large\color{white}, coltitle=white,
  title=#2, colbacktitle=lgred, toptitle=2.5pt, bottomtitle=2.5pt,
  before skip=4pt, after skip=4pt, #1
}
\newcommand{\evimg}[2]{%
  \tikz\node[inner sep=0pt, draw=cardline, line width=0.6pt, rounded corners=2pt]
    {\includegraphics[width=#2]{img/#1}};}
\newcommand{\evphoto}[2]{%
  \tikz\node[inner sep=0pt, draw=cardline, line width=0.6pt, rounded corners=2pt]
    {\includegraphics[width=\linewidth, height=#2, keepaspectratio]{img/#1}};}

\begin{document}

% ===================== CABECALHO =====================
\begin{tikzpicture}
  \fill[lgred] (0,0) rectangle (\textwidth,2.05);
  \node[anchor=west, text=white] at (0.5,1.34)
    {\lghead\fontsize{30}{34}\selectfont Evidências do Sprint 02\,{\lghead\fontsize{20}{24}\selectfont \ \textbullet\ W36-37}};
  \node[anchor=west, text=white!92] at (0.55,0.55)
    {\large Equipe 3 \;\textbullet\; Automação do Relatório de Scrap (IF Cost) \& Assistente de IA \;\textbullet\; Produção (LG)};
  \node[anchor=east] at (\textwidth-0.45,1.02) {\includegraphics[height=1.5cm]{img/logo-bw.png}};
\end{tikzpicture}\\[2pt]
{\footnotesize\color{muted} AX Academy \textbullet\ Digital Transformation \textbullet\ Programa de Equipes de Projetos \hfill Sprint 02 \textbullet\ W36-37 \textbullet\ 05 a 18 de setembro de 2026 \textbullet\ Desenvolvimento de Agentes ReAct \& Homologação}\\[6pt]

% ===================== O QUE MUDOU NESTA SPRINT (altura fixa) =====================
\begin{tcolorbox}[enhanced, colback=tintbg, colframe=lgred, boxrule=0pt,
  leftrule=3pt, arc=1pt, left=11pt, right=11pt, top=5pt, bottom=5pt,
  height=3.2cm, valign=center, before skip=0pt, after skip=7pt]
{\hd{\large O que mudou nesta sprint (Semana 8 de 11)}}\\[1pt]
{\small A frente de Inteligência Artificial do Hanaro saiu da fase de prospecção e opera com \textbf{agentes ReAct sobre o LangGraph}. O fluxo do Assistente foi apresentado e \textbf{aprovado visualmente pelo Gerente da Produção} (Leandro Corazza) em fase de PoC (o agente em si ainda não está em produção definitiva): o modelo lê dados graves do GERP (como perdas de US\$ 6.400,89), consulta a \textbf{memória de análises históricas} e sugere hipóteses de causa, contenções e tarefas. O analista atua como revisor \textit{human-in-the-loop}. O marco \textbf{M2 (03/09)} foi concluído com essa validação de PoC. \textbf{Ressalva mantida:} a automação via RPA segue incompleta pois as contas ainda não têm autorização de acesso ao GERP.}
\end{tcolorbox}

% ===================== LINHA A: 3 EVIDÊNCIAS =====================
\noindent\begin{minipage}[t]{0.335\textwidth}
\begin{card}[equal height group=evid]{Evidência 1 \textbullet\ Assistente ReAct (LangGraph)}
{\footnotesize
\textbf{IMPLEMENTADO} \ Orquestração de agentes autônomos multi-passo:\vspace{2pt}
\begin{itemize}
  \item Captura ocorrências graves de scrap a partir da ingestão do GERP (ex: item Display SOC de \textbf{US\$ 6.400,89}).
  \item Formulação automática de título executivo, categorização do tipo de defeito e descrição contextualizada.
  \item O analista avalia cada bloco individualmente (\textit{Aceitar, Editar, Rejeitar}) ou aceita a proposta completa.
  \item Fim do preenchimento manual repetitivo: esforço analítico redirecionado à decisão.
\end{itemize}}
\vspace{4pt}
{\centering\evphoto{assistente_proposta.jpg}{4.0cm}\par}
\vspace{3pt}
{\scriptsize\color{lgred}\textbf{Foco: LangGraph \textbullet\ Raciocínio ReAct \textbullet\ Sugestão de causa e contenção}}
\end{card}
\end{minipage}\hfill
\begin{minipage}[t]{0.335\textwidth}
\begin{card}[equal height group=evid]{Evidência 2 \textbullet\ Memória RAG \& Casos Parecidos}
{\footnotesize
\textbf{IMPLEMENTADO} \ Recuperação semântica em base vetorial histórica:\vspace{2pt}
\begin{itemize}
  \item Consulta à memória de análises aprovadas: o agente busca como ocorrências semelhantes de linhas e itens foram resolvidas.
  \item Respostas fundamentadas com rastreabilidade explícita das fontes (ex: casos históricos \textit{PA-ECC91611} e \textit{PA-9BD2AB95}).
  \item Interface de chat integrada (aba \textbf{Consulta}) permitindo perguntas operacionais com histórico transparente.
  \item Memória dinâmica: novas análises aprovadas pelo gestor alimentam o repositório de conhecimento da planta.
\end{itemize}}
\vspace{4pt}
{\centering\evphoto{assistente_consulta.jpg}{4.0cm}\par}
\vspace{3pt}
{\scriptsize\color{lgred}\textbf{Foco: RAG vetorial \textbullet\ Rastreabilidade de fontes \textbullet\ Aprendizado contínuo}}
\end{card}
\end{minipage}\hfill
\begin{minipage}[t]{0.315\textwidth}
\begin{card}[equal height group=evid]{Evidência 3 \textbullet\ Validação e Planos de Ação}
{\footnotesize
\textbf{PoC HOMOLOGADA} \ Revisão e geração de tarefas no relatório:\vspace{2pt}
\begin{itemize}
  \item O analista revisa hipóteses de causa e contenções, e valida tarefas corretivas sugeridas com 1 clique.
  \item Geração automática do relatório formal com plano de ação estruturado para mitigar o scrap na origem.
  \item Apresentação realizada e \textbf{aprovada visualmente pela Produção} (validação de PoC com \textbf{Leandro Corazza}; agentes em pré-produção).
  \item Marco \textbf{M2 (03/09)} consolidado; avanço para a frente de testes e homologação (M3 em 02/10).
\end{itemize}}
\vspace{4pt}
{\centering\evphoto{assistente_tarefas.jpg}{4.0cm}\par}
\vspace{3pt}
{\scriptsize\color{lgred}\textbf{Foco: Human-in-the-Loop \textbullet\ Tarefas do relatório \textbullet\ PoC Aprovada}}
\end{card}
\end{minipage}

\vspace{9pt}

% ===================== LINHA B =====================
\noindent\begin{minipage}[t]{0.38\textwidth}
\begin{card}[equal height group=fim]{A Equipe e os Resultados}
{\footnotesize
\textbf{Equipe 3:} \textbf{André Filipe Aloise} (Líder e P.O.), \textbf{Messyas Gois França} (Dev), \textbf{Gustavo Almada} (Dev) e \textbf{Teodório Neto} (Dev). Orientação: \textbf{Prof. Anderson Fontoura}.\vspace{5pt}

{\normalsize\lghead\color{lgred}Resultados e Impedimentos da Sprint 02}\vspace{1pt}
\begin{itemize}
  \item \textbf{Agentes ReAct validados visualmente:} PoC e fluxo aprovados pelo Gerente de Produção (pré-produção).
  \item \textbf{Tempo de análise reduzido em ~80\%:} Analista atua como revisor de propostas fundamentadas.
  \item \textbf{Planos de ação estruturados:} Tarefas vinculadas aos relatórios para erradicar o refugo na origem.
  \item \textbf{\color{lgred}Impedimento RPA (0\%):} Incompleto pois as contas ainda não têm autorização de acesso ao GERP.
  \item \textbf{Marco M2 concluído:} PoC validada e base pronta para os testes de sistema da Semana 9.
\end{itemize}}
\end{card}
\end{minipage}\hfill
\begin{minipage}[t]{0.37\textwidth}
\begin{card}[equal height group=fim]{Contatos e Próximos Passos}
{\footnotesize
\begin{minipage}[t]{0.5\textwidth}
{\color{lgred}\textbf{Academia (AX Academy / IFAM)}}\\[2pt]
\textbf{André Aloise}, Líder e PO\\ \faEnvelope\ aaloise@gmail.com\\[3pt]
\textbf{Prof. Anderson Fontoura}, Orientador\\ \faEnvelope\ anderson.fontoura@ifam.edu.br
\end{minipage}\hfill
\begin{minipage}[t]{0.46\textwidth}
{\color{lgred}\textbf{Cliente (LG, Produção)}}\\[2pt]
\textbf{Leandro Corazza}, valida o relatório\\ \faEnvelope\ leandro.corazza@lge.com\\[2pt]
\textbf{Lucas Cunha}, valida o relatório\\ \faEnvelope\ lucas.d.cunha@lge.com
\end{minipage}\vspace{7pt}

{\normalsize\lghead\color{lgred}Próximos passos (Semanas 9 e 10)}\vspace{1pt}
\begin{itemize}
  \item Iniciar bateria de testes de sistema, E2E e testes de segurança/LGPD (frente A4 do cronograma).
  \item Refinar filtros avançados por linha e histórico de ocorrências na Base de Scrap.
  \item Preparar ambiente de staging para homologação final e sign-off (29--30/09) rumo ao Marco M3.
\end{itemize}}
\end{card}
\end{minipage}\hfill
\begin{minipage}[t]{0.23\textwidth}
\begin{card}[equal height group=fim]{Deploy em Produção}
{\centering\evphoto{assistente_pre_analises.png}{4.4cm}\par}
\vspace{4pt}
{\footnotesize \textbf{Hanaro no Ar} \textbullet\ \textit{Intranet LG}\\[2pt]
\textbf{Ambiente Ativo:} Intranet \textbullet\ Porta \textbf{4203}\\[2pt]
{\scriptsize\color{muted}Aplicação ativa; agentes em PoC}}
\end{card}
\end{minipage}

% ===================== RODAPE =====================
\begin{tikzpicture}[remember picture, overlay]
  \fill[lgred] (current page.south west) rectangle ([yshift=0.7cm]current page.south east);
  \node[anchor=west, text=white] at ([xshift=1.5cm, yshift=0.35cm]current page.south west)
    {\footnotesize\bfseries AX Academy \textbullet\ Digital Transformation};
  \node[anchor=east, text=white!90] at ([xshift=-1.5cm, yshift=0.35cm]current page.south east)
    {\footnotesize Equipe 3 \textbullet\ Evidências do Sprint 02 (W36) \textbullet\ Automação do Relatório de Scrap};
\end{tikzpicture}

\end{document}
